\section{Introduction}
\label{sec:intro}

Most of the text a language model is trained on was written with references to other
text. A Wikipedia article on fluid dynamics links to \emph{Hydraulics} and
\emph{Archimedes' screw}; a Python file imports \texttt{chess/board.py}; a paper cites
the work it builds on. These references are not incidental. They are the author's own
statement of which other documents are needed to understand this one, and they are
recoverable from the text at essentially no cost. Standard pretraining discards them:
the corpus is shuffled into a flat token stream, documents that reference each other
land in different sequences, and even when two related documents happen to share a
window, the model is either walled off from the neighbour (document-isolated attention)
or allowed to read it indiscriminately (naive concatenation), with no notion of
\emph{which} document a given position is entitled to consult.

\method treats a corpus as a \emph{text-attributed graph}---documents are nodes, links
are directed edges---and makes the edges a first-class part of attention. Two
ingredients. First, a graph traversal packs topologically close documents into each
$32$k-token training sequence, ordered so that link targets precede the documents that
link to them. Second, a block-sparse attention mask keeps every document causally
isolated \emph{except} along explicit edges: when document $A$ links to document $B$,
the tokens of $A$ from the link position onward may attend into all of $B$. Nothing
else crosses a boundary and causality is never relaxed. The same link detector and the
same mask then run at inference: when the model emits a link that resolves to a corpus
document, that document is inserted into the context before the linking one, and the
grant opens into it exactly as it did in training.

The property we care about is the last one. Because the grant opens at the link, the
model is trained, through the ordinary next-token objective and nothing else, that
\emph{writing a reference is what makes a document readable}. A model with this
inductive bias can in principle fetch from its corpus during generation by generating a
citation---no retriever, no fusion module, no instruction tuning or reinforcement
learning to teach tool use---and the train--inference mirror is a single implementation
rather than an analogy. We regard native corpus-referencing under maintained causality
as the central promise of the design. Measuring it directly requires models large
enough to generate coherently, which the models in this paper are not; what we
establish here is the prerequisite: that the mask is learned, that the learning is
attributable to the link structure rather than to extra attention, and how the benefit
depends on the density of the graph.

Three neighbouring literatures do not occupy this point (\Cref{sec:related}).
Retrieval-augmented models use links or similarity at inference but bolt retrieval on
through a separate retriever and fusion path, with no pretraining-time analogue.
Structured packing (in-context pretraining, repository-level file ordering) puts
related documents in one window but under ordinary attention, which the packing
literature itself shows is harmful across unrelated boundaries. Graph-aware
pretraining (LinkBERT, CDLM, text-attributed-graph encoders) trains on structure, but
in encoders with symmetric attention over document pairs or through message passing,
not as a directed, link-position-gated grant inside a decoder-only language model. The
concat literature also establishes what we do \emph{not} claim: co-locating related
documents under plain attention buys most of the general long-context gain, and a
merged-corpus, general-capability story is not ours.

\paragraph{Contributions.}
\begin{itemize}
  \item A pretraining framework that packs documents by graph traversal and grants
        \emph{link-position-gated}, causality-preserving cross-document attention along
        explicit edges, with custom block-sparse kernels that make it tractable at $32$k
        context (\Cref{sec:method}, \Cref{app:kernels}).
  \item Matched-compute controls (\concatlink, \concat) that attribute the
        cross-document gain to the linking inductive bias rather than to added attention
        or to co-location of linked content (\Cref{sec:results}).
  \item Evidence on two edge modalities---Wikipedia hyperlinks and code imports across
        several languages---that link-gated attention lowers the likelihood of
        link-dependent completions relative to the same tokens scored without the
        target, while leaving single-document benchmarks unchanged.
  \item A link-density dose--response: holding corpus and packing fixed and varying only
        the fraction of edges the mask may use, the cross-document benefit on code scales
        near-linearly with the number of grants the mask realizes, is captured mostly by
        a minority of edges at training time, and is dominated by the density available
        at \emph{inference}---to the point that a model trained without any
        cross-document grants exploits them at evaluation nearly as well as one trained
        with them.
  \item A density-aware batch scheduler that removes the DDP rank imbalance induced by
        variable mask sparsity (\Cref{sec:method:density}), and a link-following
        generation loop that resolves generated links against the corpus and inserts the
        target under the training mask (\Cref{sec:method:inference}).
\end{itemize}
